\begin{abstract}
Dementia with Lewy bodies (DLB) has no dedicated biomarker-confirmed cohort in the Alzheimer's Disease Neuroimaging Initiative (ADNI), whose diagnostic labels are restricted to cognitively normal, mild cognitive impairment, and dementia. This project instead uses the Amprion CSF alpha-synuclein seed amplification assay (SAA) result as a synucleinopathy-positive proxy label, and aims to characterize regional glucose-metabolism signatures of that biological state using FDG-PET radiomics, with structural MRI and CSF/blood markers as supporting modalities rather than independent classification targets. This report reviews recent literature relevant to five pipeline-design decisions that remain open before implementation: (1) whether to reuse ADNI's own cross-scanner-harmonized FDG-PET processing pipeline instead of recoregistering raw images, (2) whether to reuse ADNI's existing FreeSurfer parcellations as the region-of-interest source instead of building a new segmentation, (3) which PyRadiomics feature classes to extract, (4) whether the current plain-sum combination of dynamic PET frames needs decay correction, and (5) how to avoid cross-validation leakage given a cohort on the order of several hundred subjects. The review finds direct precedent for the ADNI-plus-SAA cohort design in a 2025 study by Mak et al.\ \cite{mak2025alphasynuclein}, strong grounds to reuse ADNI's harmonized PET pipeline \cite{jagust2024petcore,jagust2010petcore}, and clear literature evidence that the project's current frame-summation approach is methodologically nonstandard and should be corrected \cite{jagust2010petcore}. It also surfaces a concrete, previously undocumented risk of cross-validation leakage in the planned classifier stage \cite{demircioglu2021bias,demircioglu2024oversampling}.
\end{abstract}

\begin{IEEEkeywords}
Dementia with Lewy bodies, FDG-PET, radiomics, ADNI, alpha-synuclein, seed amplification assay, PyRadiomics, cross-validation.
\end{IEEEkeywords}
